%     Tokyo Debian Meeting resources
%     Copyright (C) 2012 Junichi Uekawa
%     Copyright (C) 2011, 2015, 2020 Nobuhiro Iwamatsu

%     This program is free software; you can redistribute it and/or modify
%     it under the terms of the GNU General Public License as published by
%     the Free Software Foundation; either version 2 of the License, or
%     (at your option) any later version.

%     This program is distributed in the hope that it will be useful,
%     but WITHOUT ANY WARRANTY; without even the implied warranty of
%     MERCHANTABILITY or FITNESS FOR A PARTICULAR PURPOSE.  See the
%     GNU General Public License for more details.

%     You should have received a copy of the GNU General Public License
%     along with this program; if not, write to the Free Software
%     Foundation, Inc., 51 Franklin St, Fifth Floor, Boston, MA  02110-1301 USA

%  preview (async-shell-command (concat "evince " (replace-regexp-in-string "tex$" "pdf"(buffer-file-name))) "*tex-preview*")

%%ここからヘッダ開始。

\documentclass[mingoth,a4paper]{jsarticle}
\usepackage{monthlyreport}
% 日付を定義する、毎月変わります。
\newcommand{\debmtgyear}{2021}
\newcommand{\debmtgmonth}{7}
\newcommand{\debmtgdate}{17}
% started from zero:
% (let ((year 2021) (month 5)) (+ (* (- year 2005) 12) month -1))
\newcommand{\debmtgnumber}{199}

% Needed to import pandoc-generated LaTeX documents.
% See https://stackoverflow.com/questions/40438037/tightlist-error-using-pandoc-with-markdown
\providecommand{\tightlist}{%
  \setlength{\itemsep}{0pt}\setlength{\parskip}{0pt}}

\begin{document}

\begin{titlepage}
\thispagestyle{empty}
% タイトルページ:編集必要な部分は最初のマクロに飛ばすこと

\vspace*{-2cm}
第\debmtgnumber{}回 東京エリア Debian 勉強会資料\\
\hspace*{-2cm}
\includegraphics{image-assets/dotdeb.pdf}\\
\hfill{}\debmtgyear{}年\debmtgmonth{}月\debmtgdate{}日

% ここはアップデートすること
% 全角文字にしないとフォントのサイズが合わないので注意
\rotatebox{10}{\fontsize{30}{30} {\gt ＯＳＳＧａｔｅ＆ｐｉｕｐａｒｔｓ}}\\

\vspace*{-2cm}
\hfill{}\includegraphics[height=6cm]{image-assets/openlogo-nd.eps}
\end{titlepage}

\newpage

\begin{minipage}[b]{0.2\hsize}
 \definecolor{titleback}{gray}{0.9}
 \colorbox{titleback}{\rotatebox{90}{\fontsize{80}{80} {\gt デビアン勉強会} }}
\end{minipage}
\begin{minipage}[b]{0.8\hsize}
\hrule
\vspace{2mm}
\hrule
\begin{multicols}{2}
\tableofcontents
\end{multicols}
\vspace{2mm}
\hrule
\end{minipage}

\dancersection{最近のDebian関連のミーティング報告}{杉本　典充}

\subsection{2021年6月度 東京エリア・関西合同Debian勉強会}

2021 年 6 月 19 日(土)に東京エリア Debian 勉強会と関西 Debian 勉強会の合同で
オンラインによる Debian 勉強会を開催しました。参加者は 6 名でした。

勉強会では、「bullseyeのリリースノート日本語版の翻訳レビュー」を行いました。参加者からは様々な改善提案をいただきました。以下URLのファイルにレビュー記録をまとめています。

\url{https://tokyodebian-team.pages.debian.net/2021-06_tokyodebian_bof.txt}


\dancersection{事前課題}{杉本　典充}

今回の事前課題は以下です。

\begin{enumerate}
  \item OSS Gateという取り組みを知っていますか。
  \item piupartsを知っていますか。
\end{enumerate}

%この課題に対して提出いただいた内容は以下です。

\begin{multicols}{2}
{\small
  \begin{prework}{ ����(yy\_y\_ja\_jp) }
���ޤ���������ǤϤʤ��Ǥ���... �긵�Υ�åץȥåפǤ� /boot �ˤ� ext2
 ����¾�ˤ� ext3���Ƕ�ȤäƤ���ǥ����ȥåפǤ� /boot �ˤ� ext3 ����¾
 �ˤ� LVM ��� ext4 ��ȤäƤ��ޤ���
\end{prework}

\begin{prework}{ �����ϥ� }
�ǥե���Ȥ�ext3����
�����ǥե���ȤΤޤޡ�
(NTFS���VM���᡼�����ext3�⤢�뤱��)
\end{prework}

\begin{prework}{ yos.takahashi }
ext3/4���˻ȤäƤޤ���ext3�Υǡ�������١����ˤĤ�������Linux2011ǯ1���˼�ɮ���ޤ�����
\end{prework}

\begin{prework}{ MATOHARA }
����inode �ϳ���������Ƥ���inode ��ưŪ�˳�����Ƥ���XFS �����򤹤뤳
 �Ȥ�¿���Ǥ���NILFS �Ͼ�����Ƥߤ��ΤǤ�����mount ���˰ʲ��Τ褦�ʥ��
 ���������ФƤޤ��ݤ��ʤȻפ��ޤ�����
\begin{commandline}
$ sudo mount /dev/sdb1 /mnt
mount.nilfs2: WARNING! - The NILFS on-disk format may change at any time.
mount.nilfs2: WARNING! - Do not place critical data on a NILFS filesystem. 
\end{commandline}
����¾NotePC �Ǥ�dm-crypt �ξ�˥ե����륷���ƥ���֤��ưŹ沽�����ꡢ
 eCryptfs �ǰŹ沽�����ꤷ�Ƥ��ޤ����񤭹��߻���CPU �򤫤ʤ���񤷤ޤ��ġ�
\end{prework}

\begin{prework}{ ��ޤ� }
��ext2-$>$reiserfs-$>$jfs-$>$xfs-$>$reiserfs-$>$ext3�ȻȤäƤ��ޤ�����

����:
\begin{itemize}
 \item reiserfs: ���ե������¿���ե�����Υ������������ӥ��Ӥ��Ƥ��ɤ���
       �����������դȤ����ⵤ�������θ�μ�žȬ�ݥ����ɤˡ�����
 \item jfs: ����v1.0��̾��ä�IBM���ꥨ�ʥ��������ˤ�xfs��ƨ��
 \item xfs: fsck==true�˴�ư����������ǯ�Ȥä���ΤΥޥ�����Ĵ����0byte
       �ե��������������Ѥ���줺ƨ����
\end{itemize}
������reiser4��Ķ���Ԥ��뤦���ˤ��줬�����ʤäơ����ext3�˸��경��htree�����ä������⤦Ŵ�Ĥʤ鲿�Ǥ⤤���Ǥ��������Ȥ����Ĥ�nilfs�ʤɤ˼��Ф��Ƥ��ޤ���ext3�������noatime���٤Ǥ����������aufs���碌�Ƽ�ʬ���Ȥ��Ȥ� *strap �Ķ��򥯥����˥󥰤��Ƥ��ޤ����¸����ƥ��Ȥ������Ǥ�����Ǥ���USB�����ư�Ǥ�ͭ�ѡ�

���LVM�ǤϤʤ�MD��Ȥäƾ�Ĺ���ܥХå����åפ򤷤Ƥ��ޤ���
 MD(sda,sdb,sdc)�ǹ��ۤ����̾��MD(����)�Dz�ư���Хå����åפλ���
 attach/detach�򤹤롣�֥��å���٥�ʤΤǥꥫ�Х��FSǤ���Ǥ��������̤�
 �������¾����ˡ���ʤ�������
\end{prework}

\begin{prework}{ henrich }
�ȤäƤ���֤�NTFS��Ĺ���󤸤�ʤ��Ǥ����͡����졣
�����Ρ����Ѥο������ǥ�������ext4�ǥե����ޥåȤ��ޤ����������㤤��Ƚ��ʤ��Ȥ����������Ƥ��ޤ���
\end{prework}

\begin{prework}{ emasaka }
�Ĥ뤷��FS��ȤäƤޤ�
\end{prework}

\begin{prework}{ �ܾ� }
ext3����Ѥ��Ƥ��ޤ����ä��Ѥ�ä����ȤϤ��Ƥ��ޤ���
FS����ʤ��Ǥ������Ƕ�Lenny��2TB��HDD��Ȥä���parted�äƤλȤ��ƶä��ޤ�����
\end{prework}

\begin{prework}{ ����@������ }
���������Ū�˳��Ѥ��Ƥ���ե����륷���ƥ��ReiserFS�Ǥ���
�Ż��ǻȤäƤ�Ķ���ext3�Ǥ�����ext3���Ÿ������ǥ��㡼�ʥ뤬
����Ʋ��Ǥ������ηи��ʸŤ������ͥ�Ǥ���...�ˤ����ꡢ
���ޤ꿮�Ѥ��Ƥ��ޤ���
�����ReiserFS�Ķ��ǤϤ���ޤǤνꤤ���ʤ��Ÿ����ڤä��ꡢ
�Ƥ�HDD�����줫�����ꤷ�Ƥ��ﳲ����ä��и���̵���Τǡ�
��³Ū�˻ȤäƤ��ޤ���
��ǯ����ReiserFS�Υᥤ��ȯ��(Hans Reiser)�����ᤵ��Ƥ��ޤ������ƥʥ󥹤��ۤ��Ƥ��ޤ�����
�����������θ��ReiserFS��¾�γ�ȯ�Ԥˤ���³�����ݼ餵��Ƥ���Τǡ��¿����ޤ�����
\end{prework}

\begin{prework}{ nozzy123nozzy }
\begin{enumerate}
 \item LVM�ˤĤ��Ƥϡ�CentOS5.5��Ƴ�������Τ��Τޤޤ����Ѥ��Ƥޤ�����
       ���������ƥब��äƤ���Volume̾�ϥǥե���Ȥ�����ȡʼºݤˤ�
       kickstart�ˤơˤ��ѹ����ƻȤäƤޤ����ʾ㳲���Υ���١����˺��뤿
       ���
\item ext3�ˤĤ��Ƥϡ�debian-sid�����Τޤ޻��ꤷ�Ƥ����Τ򤽤Τޤ޻Ȥ�
      �Ƥ����ꤷ�ޤ���relatime, noatime ���餤�Ͼ����ɲä��Ƥߤ����ʡ���
      �ϻפäƤޤ���
\end{enumerate}
\end{prework}

\begin{prework}{ �ޤ��������ؤ� }
\begin{itemize}
 \item Debian�Ǥ��ä˶Ťä����Ȥ�����ext3��ȤäƤޤ������ۥޥ����qcow2
       ���᡼���ǥ��������ѻ��ʳ��ϡ�LVM�ϻȤäƤޤ���
 \item �����ǰ����ü�ʤΤϡ��������DHCP�������Ѥ�Armadillo-J�ǻȤäƤ�
       ��JFFS�Ǥ����ǥե���ȤΥե����०�����Ǥϥ�֡��Ȥ�����������
       �ƽ��������Ƥ��ޤ��Τǡ�RAM�ΰ�˽񤭤��ߡ��Ÿ��ڤäƤ�ä��ʤ�
       ���������Ǥ���Debian��udhcp�Υ������ѥå���������ӥ�ɤ��ƻȤäƤޤ���
       \footnote{\url{http://d.hatena.ne.jp/mkouhei/20080601/1212330630}}
 \item ��Debian���ߤǡ���ʬ����ǰ��֥ۥåȤʤΤ�palm webOS�Ǥ��������Ubuntu��
       �������ޥ���������Τ餷���ΤǤ�����/etc/mtab�򸫤��35�Ԥ⤢�ꡢ
       ���ʤ����֤ʹ����Ǥ��͡�
\end{itemize}
\end{prework}
}
\end{multicols}

%\dancersection{Debian Trivia Quiz}{username}
%
%Debianの昨今の話題についてのQuizです。
%
%今回の出題範囲は\url{debian-devel-announce@lists.debian.org} や
%\url{https://bits.debian.org/archives.html} や
%\url{https://micronews.debian.org/} などに投稿された内容からです。
%
%\begin{multicols}{2}
% %; whizzy-master ../debianmeetingresume201211.tex
% $B0J>e$N@_Dj$r$7$F$$$k$?$a!"$3$N%U%!%$%k$G(B M-x whizzytex $B$9$k$H!"(Bwhizzytex$B$,MxMQ$G$-$^$9!#(B
%

\santaku
{DebConf13 $B$N3+:ECO$H3+:EF|$O!)(B}
{$BF|K\!"El5~ET(B 6$B7n(B20$BF|(B}
{$B%K%+%i%0%"(B $B%^%J%0%"(B 7$B7n(B8-14$BF|(B}
{$B%9%$%9!"%t%)!<%^%k%-%e(B 8$B7n(B11-18$BF|(B}
{3}
{$B%K%+%i%0%"$O(BDebConf12$B$N3+:ECO$G$9!#(B
DebConf13$B$O%9%$%9$N%-%c%s%WCO$G3+:E$G$9!#(B
6/20$B$O3'$5$sM=Dj$r6u$1$F$*$-$^$7$g$&!#(B}

\santaku
{$B@$3&$N(BWeb$B%5!<%P$G:G$b?M5$$N$"$k(BLinux $B%G%#%9%H%j%S%e!<%7%g%s(B(W3Techs$BD4$Y(B)$B$O!)(B}
{CentOS}
{Debian}
{Ubuntu}
{B}
{\url{http://w3techs.com/technologies/history_details/os-linux}$B$K7k2L$N%0%i%U$,$"$j$^$9!#(B
$B8=:_(B Linux $B$r;HMQ$7$F$$$k(B web $B%5!<%P$N(B 32.9\% $B$,(B Debian $B$rMxMQ$7$F$*$j!"$=$N3d9g$O8=:_$bA}2C$rB3$1$F$$$k$=$&$G$9!#(B}

\santaku
{Debian $B%+!<%M%k%A!<%`$N%a%s%P!<$G$"$j!"(Bkernel.org $B$N(B 3.2.y $B0BDjHG7ONs$N%a%s%F%J$G$b$"$k(B Ben Hutchings $B$5$s$,<!4|(B Debian $B0BDjHG$H0l=o$K=P2Y$5$l$k(B Linux $B%+!<%M%k$K(B (3.2 $B7ONs$N(B mainline $B$K$OL5$$(B) $BDI2C5!G=$,Ek:\$5$l$kM=Dj$G$"$k$H=R$Y$F$$$^$9!#(B
$BB?$/$NDI2CE@$NCf$K4^$^$l$J$$$b$N$O2?!)(B}
{PREEMPT\_RT}
{Hyper-V guest drivers$B$N6/2=(B}
{ARM64/AArch64$B%"!<%-%F%/%A%c%5%]!<%H(B}
{C}
{Hyper-V guest drivers$B$O(Bmainline kernel$B$G(B3.2$B$K$b4^$^$l$F$$$^$9$,!"$h$j2~A1$5$l$?(B3.4$B$+$i$N=$@5$,F3F~$5$l$^$9!#(B
PREEMPT\_RT$B$O%O!<%I%j%"%k%?%$%`$r<B8=$9$k$?$a$N(BPatch$B!"(B
linux-image-rt-amd64 , linux-image-rt-686-pae $B$N(Bmetapackage$B$G;HMQ$G$-$^$9!#(B
$B?7$7$$(BARM 64$B%S%C%H%"!<%-%F%/%A%c%5%]!<%H$O(Bmainline kernel 3.7$B$+$i(B}

\santaku
{Wookey$B$5$s$,%"%J%&%s%9$7$?(Balpha$BHG$N(BDebian port arm64 image$B$O!)(B}
{Debian/Ubuntu port image}
{Debian/KFreeBSD port image}
{Debian/GnuHurd port image}
{A}
{self-bootstrapp(non x86)$BBP1~$H$N$3$H$G$9!#(B\url{http://wiki.debian.org/Arm64Port}$B$G%9%F!<%?%9$,3NG'$G$-$^$9!#(B}

\santaku
{700,000$BHVL\$N%P%0$,Js9p$5$l$?F|$rEv$F$k(B700000thBugContest$B$N7k2L$,=P$^$7$?!#$=$NM=A[F|$HJs9pF|$O!)(B}
{2012/12/12$B$rM=A[$7$?(BDavidPrevot}
{$BM=A[F|(B:2013/02/04$B!"Js9pF|(B:2013/02/14}
{$BM=A[F|(B:2013/02/07$B!"Js9pF|(B:2013/02/14}
{$BM=A[F|(B:2013/02/14$B!"Js9pF|(B:2013/02/07}
{C}
{$B:G$b6a$$(B2013/02/14$B$rM=A[$7$?(BChristian Perrier$B$5$s$,Ev$F$^$7$?!#7k2L$O(B\url{http://wiki.debian.org/700000thBugContest}$B$G8x3+$5$l$F$$$^$9!#(B
$B$^$?!"(B800,000/1,000,000$BHVL\$N%P%0$,Js9p$5$l$kF|$rEv$F$k%3%s%F%9%H(B\url{http://wiki.debian.org/800000thBugContest}$B$b3+:E$5$l$F$$$^$9!#(B}

\santaku
{master.debian.org$B$,?7$7$$5!3#$K0\9T$5$l$^$7$?!#$3$l$O2?$N%5!<%P$G$7$g$&$+(B $B!)(B}
{@debian.org$B$N%a!<%k%5!<%P(B}
{$B%Q%C%1!<%8$N%^%9%?!<%5!<%P(B}
{$B%Q%C%1!<%8$N%9%]%s%5!<(B(mentor)$B$rC5$9%5!<%P(B}
{A}
{$B8E$$%5!<%P$O%G%#%9%/>c32Ey$,$"$C$?$N$G!"<wL?$HH=CG$5$l!"%G!<%?$,B;<:$9$kA0$K?7$7$$%5!<%P$K0\9T$5$l$^$7$?!#(Bftp-master.debian.org$B$O(BDebian$B$N(B official package $B%j%]%8%H%j$G$9!#%Q%C%1!<%8$N%9%]%s%5!<(B(mentor)$B$rC5$9$N$O(Bmentors.debian.net$B!#(B }

\santaku
{pbuilder$B$K(Bclang support$B$,DI2C$5$l$^$7$?!#C/$,=q$$$?%Q%C%A$G$7$g$&$+!)(B}
{Sylvestre Ledru}
{Junichi Uekawa}
{Hideki Yamane}
{C}
{Debian$B$N(BClang$B%5%]!<%H$OCe!9$H?J$s$G$$$^$9!#(B}

\santaku
{DPN - 2013$BG/(B3$B7n(B4$BF|9f$K<h$j>e$2$i$l$?F|K\$N%$%Y%s%H$O(B}
{Open Source Conference 2013 Tokyo/Spring}
{Open Source Conference 2013 Hamamatu}
{Open Source Conference 2013 Tokushima}
{A}
{\url{http://henrich-on-debian.blogspot.jp/2013/02/open-source-conference-2013-tokyospring.html} $B>\:Y$O8e$[$I!#(B}


%\end{multicols}


% % (query-replace-regexp "<.*?>" "")
% % (query-replace-regexp "^[    ]\+" "")

%-------------------------------------------------------------------------------
%\dancersection{title}{username}
%-------------------------------------------------------------------------------


%-------------------------------------------------------------------------------
\dancersection{piuparts について調べてみた}{杉本 典充}
%-------------------------------------------------------------------------------

\subsection{はじめに}

Debian パッケージの作り方を学ぶことは多かったのですが、作ったパッケージをテストする方法の
知識がありませんでした。

今回は作成した Debian パッケージのインストールテストを行うツールである
piuparts について調べてみました。


\subsection{piuparts と piuparts.debian.org}

piuparts とは、Debian パッケージがインストール、アップグレード、アンインストールを
正しく処理することをテストするツールです\footnote{\url{https://packages.debian.org/ja/sid/piuparts}}。
2021 年 7 月 11 日現在、Debian unstable において piuparts-1.1.3 の Debian パッケージを配布しています\footnote{Debian 10 buster では python2.7 で動くツールでしたが、Debian 11 bullseye では python3 に対応しました。}。


piuparts パッケージの説明には「chroot 環境に最小の Debian をインストールした環境を作り、その chroot 環境でパッケージのインストールとアップグレードとアンインストールを行い、インストール前とアンインストール後のディレクトリの状態を比較して、パッケージの処理に問題があるか判定する」と記述があります。


Debian デベロッパー レファレンスの Appendix「1 Debian メンテナツールの概要」に piuparts が紹介されており\footnote{\url{https://www.debian.org/doc/manuals/developers-reference/tools.ja.html\#piuparts}}、Debian パッケージを扱う場合には piuparts を知っておくとよいと記載されています。


また、Debian Project ではパッケージを自動でテストするシステム \url{https://piuparts.debian.org/} を運用しています。
このサイトでは日々アップロードされる Debian パッケージを piuparts を使ってテストを実行し続けています。
テスト結果は \url{https://piuparts.debian.org/} の web サイトを見ることができます。


\subsection{piuparts を使ったパッケージのテスト}

\subsubsection{テストする hello パッケージのビルド}

ここでは、unstableにある hello パッケージ\url{https://packages.debian.org/sid/hello} を参考にしてみます。

まずは hello の ソースパッケージをダウンロードして、debian パッケージを手元でビルドしてみます。

\begin{commandline}
# apt-get update
# apt-get install build-essential
# apt-get build-dep hello
# apt-get source hello
# cd hello-2.10/
# dch
  -> バージョンを1つ上げます
# debuild -us -us
# ls -l .. | grep deb
-rw-r--r--  1 norimitu norimitu  36356  7月 11 04:53 hello-dbgsym_2.10-2.1_amd64.deb
-rw-r--r--  1 norimitu norimitu   6224  7月 11 04:53 hello_2.10-2.1.debian.tar.xz
-rw-r--r--  1 norimitu norimitu  56488  7月 11 04:53 hello_2.10-2.1_amd64.deb
\end{commandline}

hello\_2.10-2.1\_amd64.deb パッケージができました。
これ以降の記事では、piuparts コマンドを使ってこの Debian パッケージをテストしてみます。


\subsubsection{piuparts の実行とログ}

piuparts コマンドを使うには、piuparts パッケージをインストールします。

\begin{commandline}
# apt-get install piuparts
# which piuparts
/usr/sbin/piuparts
\end{commandline}

piuparts コマンドを実行して hello パッケージのテストを実行します。
"-l" オプションは、piupartsを実行したときに出力する標準出力をログファイルに
追記するオプションです。

\begin{commandline}
# piuparts hello_2.10-2.1_amd64.deb -l /tmp/piuparts_hello.log
\end{commandline}

最初はテストするパッケージのメタ情報を dpkg --info コマンドで出力します。

\begin{commandline}
0m0.0s INFO: ------------------------------------------------------------------------------
0m0.0s INFO: To quickly glance what went wrong, scroll down to the bottom of this logfile.
0m0.0s INFO: FAQ available at https://wiki.debian.org/piuparts/FAQ
0m0.0s INFO: The FAQ also explains how to contact us in case you think piuparts is wrong.
0m0.0s INFO: ------------------------------------------------------------------------------
0m0.0s INFO: piuparts version 1.1.3 starting up.
0m0.0s INFO: Command line arguments: /usr/sbin/piuparts hello_2.10-2.1_amd64.deb -l /tmp/piuparts.log
0m0.0s INFO: Running on: Linux sid1 4.19.0-17-amd64 #1 SMP Debian 4.19.194-1 (2021-06-10) x86_64
0m0.0s DEBUG: Starting command: ['dpkg', '--info', 'hello_2.10-2.1_amd64.deb']
0m0.0s DUMP:
   new Debian package, version 2.0.
   size 56488 bytes: control archive=1872 bytes.
       758 bytes,    20 lines      control
      3601 bytes,    49 lines      md5sums
   Package: hello
   Version: 2.10-2.1
   Architecture: amd64
   Maintainer: Santiago Vila <sanvila@debian.org>
   Installed-Size: 280
   Depends: libc6 (>= 2.14)
   Conflicts: hello-traditional
   Breaks: hello-debhelper (<< 2.9)
   Replaces: hello-debhelper (<< 2.9), hello-traditional
   Section: devel
   Priority: optional
   Homepage: http://www.gnu.org/software/hello/
   Description: example package based on GNU hello
    The GNU hello program produces a familiar, friendly greeting.  It
    allows non-programmers to use a classic computer science tool which
    would otherwise be unavailable to them.
    .
    Seriously, though: this is an example of how to do a Debian package.
    It is the Debian version of the GNU Project's `hello world' program
    (which is itself an example for the GNU Project).
    0m0.0s DEBUG: Command ok: ['dpkg', '--info', 'hello_2.10-2.1_amd64.deb']
\end{commandline}

次に、debootstrap（'--no-merged-usr'オプションを指定）を実行してテストに使う chroot 環境を作ります。
このとき、/proc や /dev などを bind mountしています。

\begin{commandline}
0m0.0s DEBUG: Created temporary directory /tmp/tmphcw125x_
0m0.0s DEBUG: Setting up minimal chroot for sid at /tmp/tmphcw125x_.
0m0.0s DEBUG: Starting command: ['debootstrap', '--variant=minbase',
'--keyring=/usr/share/keyrings/debian-archive-keyring.gpg', '--include=eatmydata',
'--no-merged-usr', '--components=main', 'sid', '/tmp/tmphcw125x_', 'http://deb.debian.org/debian']
  (debootstrapの標準出力は省略)
0m41.2s DEBUG: Command ok: ['debootstrap', '--variant=minbase',
'--keyring=/usr/share/keyrings/debian-archive-keyring.gpg', '--include=eatmydata', '--no-merged-usr',
'--components=main', 'sid', '/tmp/tmphcw125x_', 'http://deb.debian.org/debian']
0m41.2s DEBUG: Starting command: ['mount', '-t', 'proc', 'proc', '/tmp/tmphcw125x_/proc']
0m41.2s DEBUG: Command ok: ['mount', '-t', 'proc', 'proc', '/tmp/tmphcw125x_/proc']
0m41.2s DEBUG: Starting command: ['mount', '-t', 'devpts',
'-o', 'newinstance,noexec,nosuid,gid=5,mode=0620,ptmxmode=0666', 'devpts', '/tmp/tmphcw125x_/dev/pts']
0m41.2s DEBUG: Command ok: ['mount', '-t', 'devpts',
'-o', 'newinstance,noexec,nosuid,gid=5,mode=0620,ptmxmode=0666', 'devpts', '/tmp/tmphcw125x_/dev/pts']
0m41.2s DEBUG: Starting command: ['mount', '-o', 'bind', '/tmp/tmphcw125x_/dev/pts/ptmx',
'/tmp/tmphcw125x_/dev/ptmx']
0m41.2s DEBUG: Command ok: ['mount', '-o', 'bind', '/tmp/tmphcw125x_/dev/pts/ptmx', '/tmp/tmphcw125x_/dev/ptmx']
0m41.2s DEBUG: Starting command: ['mount', '-o', 'bind', '/dev/pts/4', '/tmp/tmphcw125x_/dev/console']
0m41.2s DEBUG: Command ok: ['mount', '-o', 'bind', '/dev/pts/4', '/tmp/tmphcw125x_/dev/console']
0m41.2s DEBUG: Starting command: ['mount', '-t', 'tmpfs', '-o', 'size=65536k', 'tmpfs', '/tmp/tmphcw125x_/dev/shm']
0m41.2s DEBUG: Command ok: ['mount', '-t', 'tmpfs', '-o', 'size=65536k', 'tmpfs', '/tmp/tmphcw125x_/dev/shm']
0m41.2s DEBUG: sources.list:
  deb http://deb.debian.org/debian sid main
0m41.2s DEBUG: Created policy-rc.d and chmodded it.
0m41.2s DEBUG: Created resolv.conf.
\end{commandline}

debootstrap で作成したディレクトリに chroot で入り、apt-get update を実行して試験対象のリポジトリの
パッケージ情報を取得します。また、このタイミングで chroot のテスト環境のディレクトリ状態を保存しています。

\begin{commandline}
0m41.2s DEBUG: Starting command: ['chroot', '/tmp/tmphcw125x_', 'eatmydata', 'apt-get', 'update']
0m44.4s DUMP:
  Hit:1 http://deb.debian.org/debian sid InRelease
  Get:2 http://deb.debian.org/debian sid/main Translation-en [6509 kB]
  Fetched 6509 kB in 2s (2650 kB/s)
  Reading package lists...
0m44.4s DEBUG: Command ok: ['chroot', '/tmp/tmphcw125x_', 'eatmydata', 'apt-get', 'update']
0m44.4s DEBUG: Starting command: ['chroot', '/tmp/tmphcw125x_', 'eatmydata', 'sh', '-c', 'apt-cache dumpavail | md5sum']
0m44.5s DUMP:
  c0117488ce119ec145559966f8a0cef8  -
0m44.5s DEBUG: Command ok: ['chroot', '/tmp/tmphcw125x_', 'eatmydata', 'sh', '-c', 'apt-cache dumpavail | md5sum']
0m44.5s DEBUG: Starting command: ['chroot', '/tmp/tmphcw125x_', 'eatmydata', 'apt-get', 'clean']
0m44.5s DEBUG: Command ok: ['chroot', '/tmp/tmphcw125x_', 'eatmydata', 'apt-get', 'clean']
0m44.5s DEBUG: Recording chroot state
\end{commandline}

dpkg-query、dpkg-divert --listを実行してパッケージの情報や中身を調べています。

\begin{commandline}
0m44.6s DEBUG: Starting command: ['chroot', '/tmp/tmphcw125x_', 'eatmydata', 'dpkg-query', '-W',
'-f', '${Status}\\t${binary:Package}\\t${Package}\\t${Version}\\n']
　（dpkg-queryの結果は省略）
0m44.6s DEBUG: Command ok: ['chroot', '/tmp/tmphcw125x_', 'eatmydata', 'dpkg-query', '-W',
'-f', '${Status}\\t${binary:Package}\\t${Package}\\t${Version}\\n']
0m44.6s DEBUG: Starting command: ['chroot', '/tmp/tmphcw125x_', 'eatmydata', 'dpkg-divert', '--list']
0m44.6s DUMP:
  diversion of /usr/share/man/man1/sh.1.gz to /usr/share/man/man1/sh.distrib.1.gz by dash
  diversion of /bin/sh to /bin/sh.distrib by dash
0m44.6s DEBUG: Command ok: ['chroot', '/tmp/tmphcw125x_', 'eatmydata', 'dpkg-divert', '--list']
0m44.6s INFO: apt-cache does not know about any of the requested packages
0m44.6s DEBUG: Starting command: ['lsof', '-w', '+D', '/tmp/tmphcw125x_']
0m44.8s DEBUG: Command failed (status=1), but ignoring error: ['lsof', '-w', '+D', '/tmp/tmphcw125x_']
0m45.3s DEBUG: No broken symlinks as far as we can find.
0m45.3s DEBUG: Starting command: ['lsof', '-w', '+D', '/tmp/tmphcw125x_']
0m45.4s DEBUG: Command failed (status=1), but ignoring error: ['lsof', '-w', '+D', '/tmp/tmphcw125x_']
0m45.9s DEBUG: No broken symlinks as far as we can find.
0m45.9s DEBUG: Starting command: ['chroot', '/tmp/tmphcw125x_', 'eatmydata', 'dpkg-query', '-f', '${Version}\n', '-W', 'apt']
0m45.9s DUMP:
  2.2.4
0m45.9s DEBUG: Command ok: ['chroot', '/tmp/tmphcw125x_', 'eatmydata', 'dpkg-query', '-f', '${Version}\n', '-W', 'apt']
\end{commandline}

piuparts コマンドの引数に指定した Debian パッケージを chroot のテスト環境内にコピーし、
apt-get install コマンドを実行してコピーした Debian パッケージをインストールします。

\begin{commandline}
0m45.9s DEBUG: Copying hello_2.10-2.1_amd64.deb to /tmp/tmphcw125x_/tmp
0m45.9s DEBUG: Starting command: ['chroot', '/tmp/tmphcw125x_', 'eatmydata', 'apt-get', '-y',
'--allow-downgrades', 'install', './tmp/hello_2.10-2.1_amd64.deb']
0m47.1s DUMP:
  Reading package lists...
  Building dependency tree...
  The following NEW packages will be installed:
    hello
  0 upgraded, 1 newly installed, 0 to remove and 0 not upgraded.
  Need to get 0 B/56.5 kB of archives.
  After this operation, 287 kB of additional disk space will be used.
  Get:1 /tmp/hello_2.10-2.1_amd64.deb hello amd64 2.10-2.1 [56.5 kB]
  debconf: delaying package configuration, since apt-utils is not installed
  Selecting previously unselected package hello.

  Preparing to unpack /tmp/hello_2.10-2.1_amd64.deb ...
  Unpacking hello (2.10-2.1) ...
  Setting up hello (2.10-2.1) ...
0m47.1s DEBUG: Command ok: ['chroot', '/tmp/tmphcw125x_', 'eatmydata', 'apt-get', '-y',
'--allow-downgrades', 'install', './tmp/hello_2.10-2.1_amd64.deb']
0m47.1s DEBUG: Starting command: ['chroot', '/tmp/tmphcw125x_', 'eatmydata', 'dpkg-query',
'-f', '${Package} ${Status}\n', '-W', 'hello']
0m47.1s DUMP:
  hello install ok installed
0m47.1s DEBUG: Command ok: ['chroot', '/tmp/tmphcw125x_', 'eatmydata', 'dpkg-query',
'-f', '${Package} ${Status}\n', '-W', 'hello']
0m47.1s INFO: Installation of ['./tmp/hello_2.10-2.1_amd64.deb'] ok
\end{commandline}

chroot のテスト環境内のディレクトリを調べたり、インストールしているパッケージのバージョンを調べています。

\begin{commandline}
0m47.1s DEBUG: Removing /tmp/tmphcw125x_/./tmp/hello_2.10-2.1_amd64.deb
0m47.1s DEBUG: Starting command: ['lsof', '-w', '+D', '/tmp/tmphcw125x_']
0m47.3s DEBUG: Command failed (status=1), but ignoring error: ['lsof', '-w', '+D', '/tmp/tmphcw125x_']
0m47.7s DEBUG: No broken symlinks as far as we can find.
0m48.0s DEBUG: Starting command: ['debsums', '--root', '/tmp/tmphcw125x_', '-ac', '--ignore-obsolete']
0m48.5s DEBUG: Command ok: ['debsums', '--root', '/tmp/tmphcw125x_', '-ac', '--ignore-obsolete']
0m48.5s DEBUG: Starting command: ['dpkg-query', '-f', '${Version}\n', '-W', 'adequate']
0m48.5s DUMP:
  0.15.6
0m48.5s DEBUG: Command ok: ['dpkg-query', '-f', '${Version}\n', '-W', 'adequate']
0m48.5s INFO: Running adequate version 0.15.6 now.
0m48.5s DEBUG: Starting command: ['adequate', '--root', '/tmp/tmphcw125x_', 'hello']
0m48.6s DEBUG: Command ok: ['adequate', '--root', '/tmp/tmphcw125x_', 'hello']
0m48.6s DEBUG: Starting command: ['chroot', '/tmp/tmphcw125x_', 'eatmydata', 'dpkg-query', '-W',
'-f', '${Status}\\t${binary:Package}\\t${Package}\\t${Version}\\n'
0m48.6s DEBUG: Command ok: ['chroot', '/tmp/tmphcw125x_', 'eatmydata', 'dpkg-query', '-W',
'-f', '${Status}\\t${binary:Package}\\t${Package}\\t${Version}\\n']
\end{commandline}

chroot のテスト環境内でインストールしたパッケージを apt-get remove を実行してアンインストールします。
これで Debian パッケージのインストールからアンインストールするテストは終了です。

\begin{commandline}
0m48.6s DEBUG: Starting command: ['chroot', '/tmp/tmphcw125x_', 'eatmydata', 'apt-get',
'remove', 'hello']
  （省略）
0m49.1s DEBUG: Command ok: ['chroot', '/tmp/tmphcw125x_', 'eatmydata', 'apt-get', 'remove', 'hello']
0m49.1s DEBUG: Starting command: ['chroot', '/tmp/tmphcw125x_', 'eatmydata', 'dpkg', '--purge', 'hello']
0m49.2s DUMP:
  dpkg: warning: ignoring request to remove hello which isn't installed
0m49.2s DEBUG: Command ok: ['chroot', '/tmp/tmphcw125x_', 'eatmydata', 'dpkg', '--purge', 'hello']
0m49.2s DEBUG: Starting command: ['chroot', '/tmp/tmphcw125x_', 'eatmydata', 'dpkg', '--purge', '--pending']
0m49.2s DEBUG: Command ok: ['chroot', '/tmp/tmphcw125x_', 'eatmydata', 'dpkg', '--purge', '--pending']
0m49.2s DEBUG: Starting command: ['chroot', '/tmp/tmphcw125x_', 'eatmydata', 'dpkg', '--remove', '--pending']
0m49.2s DEBUG: Command ok: ['chroot', '/tmp/tmphcw125x_', 'eatmydata', 'dpkg', '--remove', '--pending']
0m49.2s DEBUG: Starting command: ['lsof', '-w', '+D', '/tmp/tmphcw125x_']
0m49.4s DEBUG: Command failed (status=1), but ignoring error: ['lsof', '-w', '+D', '/tmp/tmphcw125x_']
0m49.8s DEBUG: No broken symlinks as far as we can find.
0m49.8s DEBUG: Starting command: ['chroot', '/tmp/tmphcw125x_', 'eatmydata', 'dpkg-divert', '--list']
0m49.8s DUMP:
  diversion of /usr/share/man/man1/sh.1.gz to /usr/share/man/man1/sh.distrib.1.gz by dash
  diversion of /bin/sh to /bin/sh.distrib by dash
0m49.8s DEBUG: Command ok: ['chroot', '/tmp/tmphcw125x_', 'eatmydata', 'dpkg-divert', '--list']
0m49.8s DEBUG: Starting command: ['chroot', '/tmp/tmphcw125x_', 'eatmydata', 'apt-get', 'clean']
0m49.9s DEBUG: Command ok: ['chroot', '/tmp/tmphcw125x_', 'eatmydata', 'apt-get', 'clean']
0m49.9s DEBUG: Recording chroot state
0m50.0s INFO: PASS: Installation and purging test.
\end{commandline}

引き続き、インストール、アップグレード、アンインストールのテストを実行します。
まずはリポジトリにある同名の hello パッケージを apt-get install でインストールします。

\begin{commandline}
0m50.0s DEBUG: Starting command: ['chroot', '/tmp/tmphcw125x_', 'eatmydata', 'apt-cache', 'show',
  '--no-all-versions', 'hello']
0m50.9s DEBUG: Command ok: ['chroot', '/tmp/tmphcw125x_', 'eatmydata', 'apt-cache', 'policy', 'hello']
0m50.9s DEBUG: Starting command: ['chroot', '/tmp/tmphcw125x_', 'eatmydata', 'apt-get', '-y', 'install', 'hello']
0m51.5s DUMP:
  Reading package lists...
  Building dependency tree...
  Reading state information...
  The following NEW packages will be installed:
    hello
    0 upgraded, 1 newly installed, 0 to remove and 0 not upgraded.
0m51.5s DEBUG: Command ok: ['chroot', '/tmp/tmphcw125x_', 'eatmydata', 'apt-get', '-y', 'install', 'hello']
0m51.5s DEBUG: Starting command: ['lsof', '-w', '+D', '/tmp/tmphcw125x_']
0m51.7s DEBUG: Command failed (status=1), but ignoring error: ['lsof', '-w', '+D', '/tmp/tmphcw125x_']
0m52.2s DEBUG: No broken symlinks as far as we can find.
0m52.2s DEBUG: Starting command: ['chroot', '/tmp/tmphcw125x_', 'eatmydata', 'dpkg-query', '-f', '${Version}\n', '-W', 'apt']
0m52.2s DUMP:
  2.2.4
0m52.2s DEBUG: Command ok: ['chroot', '/tmp/tmphcw125x_', 'eatmydata', 'dpkg-query', '-f', '${Version}\n', '-W', 'apt']
\end{commandline}

次に piuparts コマンドの引数に指定した Debian パッケージを chroot のテスト環境内にコピーし、
apt-get install コマンドを実行してコピーしたパッケージでアップグレードします。

\begin{commandline}
0m52.2s DEBUG: Copying hello_2.10-2.1_amd64.deb to /tmp/tmphcw125x_/tmp
0m52.2s DEBUG: Starting command: ['chroot', '/tmp/tmphcw125x_', 'eatmydata', 'apt-get', '-y',
'--allow-downgrades', 'install', './tmp/hello_2.10-2.1_amd64.deb']
0m52.8s DUMP:
  Reading package lists...
  Building dependency tree...
  Reading state information...
  The following packages will be upgraded:
    hello
  1 upgraded, 0 newly installed, 0 to remove and 0 not upgraded.

0m52.8s DEBUG: Command ok: ['chroot', '/tmp/tmphcw125x_', 'eatmydata', 'apt-get', '-y',
'--allow-downgrades', 'install', './tmp/hello_2.10-2.1_amd64.deb']
0m52.8s DEBUG: Starting command: ['chroot', '/tmp/tmphcw125x_', 'eatmydata', 'dpkg-query',
'-f', '${Package} ${Status}\n', '-W', 'hello']
0m52.8s DUMP:
  hello install ok installed
0m52.8s DEBUG: Command ok: ['chroot', '/tmp/tmphcw125x_', 'eatmydata', 'dpkg-query',
'-f', '${Package} ${Status}\n', '-W', 'hello']
0m52.8s INFO: Installation of ['./tmp/hello_2.10-2.1_amd64.deb'] ok
\end{commandline}

次に apt-get remove コマンドを実行してパッケージをアンインストールします。
これで Debian パッケージのインストールからアップグレード、アンインストールするテストは終了です。

\begin{commandline}
0m54.3s DEBUG: Starting command: ['chroot', '/tmp/tmphcw125x_', 'eatmydata', 'apt-get', 'remove', 'hello']
0m54.9s DUMP:
  Reading package lists...
  Building dependency tree...
  Reading state information...
  The following packages will be REMOVED:
    hello
  0 upgraded, 0 newly installed, 1 to remove and 0 not upgraded.
0m54.9s DEBUG: Command ok: ['chroot', '/tmp/tmphcw125x_', 'eatmydata', 'apt-get', 'remove', 'hello']
0m54.9s DEBUG: Starting command: ['chroot', '/tmp/tmphcw125x_', 'eatmydata', 'dpkg', '--purge', 'hello']
0m54.9s DUMP:
  dpkg: warning: ignoring request to remove hello which isn't installed
0m54.9s DEBUG: Command ok: ['chroot', '/tmp/tmphcw125x_', 'eatmydata', 'dpkg', '--purge', 'hello']
0m54.9s DEBUG: Starting command: ['chroot', '/tmp/tmphcw125x_', 'eatmydata', 'dpkg', '--purge', '--pending']
0m54.9s DEBUG: Command ok: ['chroot', '/tmp/tmphcw125x_', 'eatmydata', 'dpkg', '--purge', '--pending']
0m54.9s DEBUG: Starting command: ['chroot', '/tmp/tmphcw125x_', 'eatmydata', 'dpkg', '--remove', '--pending']
0m54.9s DEBUG: Command ok: ['chroot', '/tmp/tmphcw125x_', 'eatmydata', 'dpkg', '--remove', '--pending']
0m54.9s DEBUG: Starting command: ['lsof', '-w', '+D', '/tmp/tmphcw125x_']
0m55.1s DEBUG: Command failed (status=1), but ignoring error: ['lsof', '-w', '+D', '/tmp/tmphcw125x_']
0m55.5s DEBUG: No broken symlinks as far as we can find.
0m55.5s DEBUG: Starting command: ['chroot', '/tmp/tmphcw125x_', 'eatmydata', 'dpkg-divert', '--list']
0m55.5s DUMP:
  diversion of /usr/share/man/man1/sh.1.gz to /usr/share/man/man1/sh.distrib.1.gz by dash
  diversion of /bin/sh to /bin/sh.distrib by dash
0m55.5s DEBUG: Command ok: ['chroot', '/tmp/tmphcw125x_', 'eatmydata', 'dpkg-divert', '--list']
0m55.5s DEBUG: Starting command: ['chroot', '/tmp/tmphcw125x_', 'eatmydata', 'apt-get', 'clean']
0m55.5s DEBUG: Command ok: ['chroot', '/tmp/tmphcw125x_', 'eatmydata', 'apt-get', 'clean']
0m55.5s DEBUG: Recording chroot state
0m55.7s INFO: PASS: Installation, upgrade and purging tests.
\end{commandline}

最後に debootstrap で構築した chroot 環境をクリーンナップします。

\begin{commandline}
0m55.9s DEBUG: Starting command: ['umount', '/tmp/tmphcw125x_/dev/shm']
0m55.9s DEBUG: Command ok: ['umount', '/tmp/tmphcw125x_/dev/shm']
0m55.9s DEBUG: Starting command: ['umount', '/tmp/tmphcw125x_/dev/console']
0m55.9s DEBUG: Command ok: ['umount', '/tmp/tmphcw125x_/dev/console']
0m55.9s DEBUG: Starting command: ['umount', '/tmp/tmphcw125x_/dev/ptmx']
0m55.9s DEBUG: Command ok: ['umount', '/tmp/tmphcw125x_/dev/ptmx']
0m55.9s DEBUG: Starting command: ['umount', '/tmp/tmphcw125x_/dev/pts']
0m56.0s DEBUG: Command ok: ['umount', '/tmp/tmphcw125x_/dev/pts']
0m56.0s DEBUG: Starting command: ['umount', '/tmp/tmphcw125x_/proc']
0m56.0s DEBUG: Command ok: ['umount', '/tmp/tmphcw125x_/proc']
0m56.0s DEBUG: Starting command: ['rm', '-rf', '--one-file-system', '/tmp/tmphcw125x_']
0m56.1s DEBUG: Command ok: ['rm', '-rf', '--one-file-system', '/tmp/tmphcw125x_']
0m56.1s DEBUG: Removed directory tree at /tmp/tmphcw125x_
0m56.1s INFO: PASS: All tests.
0m56.1s INFO: piuparts run ends.
\end{commandline}


\subsection{piuparts コマンドのオプションの紹介}

\subsubsection{-d オプション：バージョンを指定してテスト}

piuparts に "-d \_distro\_" オプションを指定すると、deboostrap する環境のバージョンを
指定してテストを実行します。"-d" オプションは省略でき、省略した場合は "-d sid" を
指定したときと同じ動きをします。

\begin{commandline}
# piuparts hello_2.10-2.1_amd64.deb -d sid -l /tmp/piuparts.log
\end{commandline}

この "-d" オプションは複数を連続して指定することができます。
その場合、最初に指定した"-d" オプションのバージョンで debootstrap してからパッケージをインストールし、それ以降に指定した "-d" のバージョンへ順番に apt-get dist-upgrade を行い、最後まで dist-upgrade が完了した状態で piuparts コマンドに指定したパッケージのインストールテストを実行します。

以下の例では、buster の環境を deboostrap で構築して hello パッケージをインストールし、
bullseye へ apt-get dist-upgrade し、sid へapt-get dist-upgrade した後に、
hello\_2.10-2.1\_amd64.deb のアップグレードとアンインストールのテストを実行します。

\begin{commandline}
# piuparts hello_2.10-2.1_amd64.deb -d buster -d bullseye -d sid -l /tmp/piuparts.log
\end{commandline}


\subsubsection{インストール済みパッケージのテスト}

これまでの紹介では Debian パッケージ自体を piuparts の引数に指定してきました。
テストするパッケージは .deb ファイルではなく、"--apt"オプションとパッケージ名を指定することで
実行している PC にインストールしている Debian パッケージのテストを行うことができます。

\begin{commandline}
# piuparts hello --apt -l /tmp/piuparts.log
\end{commandline}


\subsubsection{-p オプション：pbuilderのイメージを使う}

Debian パッケージをクリーン環境でビルドするツールに pbuiler\footnote{\url{https://packages.debian.org/ja/sid/pbuilder}} があります。

以下のように pbuilder コマンドを実行して debootstrap したクリーン環境で Debian パッケージのビルドを行うことができます。

\begin{commandline}
# pbuilder create
# pbuilder update
# pbuilder build hello_2.10-2.1.dsc
\end{commandline}

pbuilder create を実行したとき、または pbuilder update を実行したときに pbuilder が利用する
deboostrap のイメージを「/var/cache/pbuilder/base.tgz」に保存します。

piuparts に "-p" オプションを指定すると、debootstrap を実行する代わりに pbuilder が保存している
「/var/cache/pbuilder/base.tgz」をテスト環境に利用します。

deboostrap を実行するとインターネットから毎回約 100 MByte 程度のファイルをダウンロードするため
"-p" オプションを使うことでネットワーク帯域やダウンロード時間を節約できます。

似たようなオプションに "--basetgz=file" があり、base.tgz ファイルを任意のパスで指定することもできます。


\subsubsection{--docker-image：dockerイメージを使ってテストする}

piuparts に "--docker-image=IMAGE" オプションを指定すると、テストに使う環境に
debootstrap で構築したイメージを使わず、dockerイメージを使ってテストを実行します。

実際試してみると、chroot コマンドを使って chroot のテスト環境内でコマンドを実行するところが
docker run、docker inspect、docker top、docker exec などを使ってテストを実行する動きに変わります。


\subsubsection{--merged-usr：/usr マージした環境でdebootstrapする}

piuparts に "--merged-usr" オプションを指定すると、debootstrap するときに "--merged-usr" オプションを指定して実行します（オプションなしの場合は "--no-merged-usr" を指定しています）。

これで /usr マージした環境で Debian パッケージのインストールテストを実行することができます。


\subsubsection{--arch：アーキテクチャを指定してdeboostrapする}

piuparts に "--arch=ARCH" オプションを指定すると、debootstrap するときに
"--arch=ARCH" を指定して実行します。

つまり、異なるアーキテクチャの chroot 環境を構築して Debian パッケージのインストールテストを
実行することができます。
ただ、テストするPCが --arch オプションに指定したアーキテクチャで動作する環境であること、
Debian パッケージも --arch オプションに指定したアーキテクチャ向けにビルドしてある必要があります。


\subsection{終わりに}

piuparts を使ってパッケージのテストを試してみました。
今回試した範囲では Debian Developer が作成したテストが合格するパッケージのみでした。
今後パッケージを自分で作成するときには様々な不具合に遭遇と思いますので、piuparts を使ってうまくテストできるようになれるとよいと思います。

なお、piuparts コマンドはネットワーク帯域を多く使うため、回線の帯域や通信量の残量に注意してご利用ください。
\footnote{\url{https://piuparts.debian.org/doc/README_server.txt}に piuparts のサーバでは 100 Mbps の帯域を簡単に使い切ってしまうから気をつけてください、と注意書きがあります。}

\subsection{参考資料}

\begin{itemize}
\item Debian Wiki - piuparts \url{https://wiki.debian.org/piuparts}
\item 岩松 信洋、「piuparts の使い方」、あんどきゅめんてっど でびあん 2010 年夏号 \\
\url{https://tokyodebian-team.pages.debian.net/pdf2010/debianmeetingresume201004.pdf}
\item Holger Levsen、DebConf19「i'm (a bit) sick of maintaining piuparts.debian.org (mostly) alone, please help」\footnote{\url{https://debconf19.debconf.org/talks/103-im-a-bit-sick-of-maintaining-piupartsdebianorg-mostly-alone-please-help/}}
\end{itemize}


% 冊子にするために、4の倍数にする必要がある。
% そのための調整
% \dancersection{メモ}{}
% \mbox{}\newpage

\vspace*{15cm}
\hrule
\vspace{2mm}
\includegraphics[width=2cm]{image-assets/openlogo-nd.eps}
\noindent \Large \bf Debian 勉強会資料\\
\noindent \normalfont \debmtgyear{}年\debmtgmonth{}月\debmtgdate{}日 \hspace{5mm}  初版第1刷発行\\
\noindent \normalfont 東京エリア Debian 勉強会 （編集・印刷・発行）\\
\hrule
\end{document}
